% GENERATED FILE. Edit canonical lessons, then run npm run generate:books.
% canonical-source-hash: fnv1a64:331859c97c75b596
% canonical-lessons: LA-C430-temo, LA-C430-modiolus, LA-C430-sufflamen, LA-C430-subiungo, LA-C430-disiungo, LA-R430-review-423-430

\chapter{Cart pole, Wheel hub, Drag shoe, To yoke, To unyoke}
\label{ch:la-cart-pole-wheel-hub-drag-shoe-to-yoke-to-unyoke}
\hlchaptermodality{430}

\begin{chapteropening}
\textbf{By the end of this chapter:} \emph{I can talk about the parts of a cart and harnessing it: the pole, the hub, the drag shoe, yoking and unyoking.}

Complete the last lesson of chapter 430: I can talk about the parts of a cart and harnessing it: the pole, the hub, the drag shoe, yoking and unyoking.
\end{chapteropening}

\section[tēmō, tēmōnis]{tēmō, tēmōnis — a cart pole, the shaft of a wagon or plough}

\label{lesson:LA-C430-temo}



\begin{quote}
Before the new one: say the Latin for a covered carriage, then the Latin for a travelling coach.
\end{quote}

\subsection*{You'll want to know: tēmō, tēmōnis}
\textbf{tēmō, tēmōnis} — \textquotedblleft{}a cart pole, the shaft of a wagon or plough\textquotedblright{}.

\begin{quote}
Agricola bovēs ad tēmōnem plaustrī iungit.
\end{quote}

\begin{quote}
\emph{The farmer harnesses the oxen to the pole of the wagon.}
\end{quote}

Third declension, masculine: \emph{tēmō, tēmōnis m.} The long beam that runs forward from the cart to the yoke.

\begin{cousinweb}
\emph{tēmō} is inherited. Lewis \& Short connect it with the root \emph{tek-} of \emph{texō} 'to weave, to build', and EtymDB lists \emph{texō} as a cousin; the exact path is not settled.
\end{cousinweb}

\begin{grammarlens}[title={What you've built}]
One more word for telling, explaining and describing at B1.
\end{grammarlens}

\subsection*{Your turn}
\begin{itemize}
\raggedright
  \item \emph{Say it:} \emph{tēmō, tēmōnis}
  \item \emph{Say it:} \emph{tēmō, tēmōnis}, once more
  \item \emph{Say it:} a sentence of your own with \emph{tēmō, tēmōnis}
  \item \emph{Recall:} say the Latin for a covered carriage, then the Latin for a travelling coach, then say \emph{tēmō, tēmōnis} again
\end{itemize}

\begin{tcolorbox}[breakable,colback=teal!4,colframe=teal!35!black,title={Before you move on}]
What does tēmō, tēmōnis mean? (\textquotedblleft{}A cart pole, the shaft of a wagon or plough\textquotedblright{}.) Say it once more.
\end{tcolorbox}
\section[modiolus, modiolī]{modiolus, modiolī — a wheel hub, the nave of a wheel; a small measure}

\label{lesson:LA-C430-modiolus}



\begin{quote}
Before the new one: say the Latin for a travelling coach, then the Latin for a cart pole.
\end{quote}

\subsection*{You'll want to know: modiolus, modiolī}
\textbf{modiolus, modiolī} — \textquotedblleft{}a wheel hub, the nave of a wheel; a small measure\textquotedblright{}.

\begin{quote}
In mediā rotā modiolus est.
\end{quote}

\begin{quote}
\emph{In the middle of the wheel is the hub.}
\end{quote}

Second declension, masculine: \emph{modiolus, modiolī m.} A little measure: the hub looks like a small barrel.

\begin{cousinweb}
\emph{modiolus} is a diminutive of \emph{modius} 'a peck measure' (Lewis \& Short), which is built on \emph{modus} 'measure' (EtymDB). The round hub of a wheel looked like a small measuring tub.
\end{cousinweb}

\begin{grammarlens}[title={What you've built}]
One more word for telling, explaining and describing at B1.
\end{grammarlens}

\subsection*{Your turn}
\begin{itemize}
\raggedright
  \item \emph{Say it:} \emph{modiolus, modiolī}
  \item \emph{Say it:} \emph{modiolus, modiolī}, once more
  \item \emph{Say it:} a sentence of your own with \emph{modiolus, modiolī}
  \item \emph{Recall:} say the Latin for a travelling coach, then the Latin for a cart pole, then say \emph{modiolus, modiolī} again
\end{itemize}

\begin{tcolorbox}[breakable,colback=teal!4,colframe=teal!35!black,title={Before you move on}]
What does modiolus, modiolī mean? (\textquotedblleft{}A wheel hub, the nave of a wheel; a small measure\textquotedblright{}.) Say it once more.
\end{tcolorbox}
\section[sufflāmen, sufflāminis]{sufflāmen, sufflāminis — a drag shoe, a brake (for a wheel)}

\label{lesson:LA-C430-sufflamen}



\begin{quote}
Before the new one: say the Latin for a cart pole, then the Latin for a wheel hub.
\end{quote}

\subsection*{You'll want to know: sufflāmen, sufflāminis}
\textbf{sufflāmen, sufflāminis} — \textquotedblleft{}a drag shoe, a brake (for a wheel)\textquotedblright{}.

\begin{quote}
In clīvō raedārius sufflāmen sub rotā pōnit.
\end{quote}

\begin{quote}
\emph{On the slope the coachman puts the drag shoe under the wheel.}
\end{quote}

Third declension, neuter: \emph{sufflāmen, sufflāminis n.} A block or chain that stops a wheel turning on a steep road. Also anything that holds a thing back.

\begin{cousinweb}
The origin of \emph{sufflāmen} is not settled. Lewis \& Short file it with \emph{sufflō} 'to blow up' without further comment, while EtymDB lists Greek \emph{phalanx} 'log, beam' as a cousin, which points to a word for a wooden block.
\end{cousinweb}

\begin{grammarlens}[title={What you've built}]
One more word for telling, explaining and describing at B1.
\end{grammarlens}

\subsection*{Your turn}
\begin{itemize}
\raggedright
  \item \emph{Say it:} \emph{sufflāmen, sufflāminis}
  \item \emph{Say it:} \emph{sufflāmen, sufflāminis}, once more
  \item \emph{Say it:} a sentence of your own with \emph{sufflāmen, sufflāminis}
  \item \emph{Recall:} say the Latin for a cart pole, then the Latin for a wheel hub, then say \emph{sufflāmen, sufflāminis} again
\end{itemize}

\begin{tcolorbox}[breakable,colback=teal!4,colframe=teal!35!black,title={Before you move on}]
What does sufflāmen, sufflāminis mean? (\textquotedblleft{}A drag shoe, a brake (for a wheel)\textquotedblright{}.) Say it once more.
\end{tcolorbox}
\section[subiungō, subiungere]{subiungō, subiungere — to yoke, to harness; to put under}

\label{lesson:LA-C430-subiungo}



\begin{quote}
Before the new one: say the Latin for a wheel hub, then the Latin for a drag shoe.
\end{quote}

\subsection*{You'll want to know: subiungō, subiungere}
\textbf{subiungō, subiungere} — \textquotedblleft{}to yoke, to harness; to put under\textquotedblright{}.

\begin{quote}
Agricola mūlōs plaustrō subiungit.
\end{quote}

\begin{quote}
\emph{The farmer harnesses the mules to the wagon.}
\end{quote}

\emph{subiungō, subiungere, subiūnxī, subiūnctum}, often with a dative: \emph{equōs currū subiungere}. English \emph{subjunctive} comes from it.

\begin{cousinweb}
\emph{subiungō} is built from \emph{sub-} 'under' and \emph{iungō} 'to join, to yoke' (EtymDB; Lewis \& Short). English \emph{subjunctive} comes from its participle \emph{subiūnctus}.
\end{cousinweb}

\begin{grammarlens}[title={What you've built}]
One more word for telling, explaining and describing at B1.
\end{grammarlens}

\subsection*{Your turn}
\begin{itemize}
\raggedright
  \item \emph{Say it:} \emph{subiungō, subiungere}
  \item \emph{Say it:} \emph{subiungō, subiungere}, once more
  \item \emph{Say it:} a sentence of your own with \emph{subiungō, subiungere}
  \item \emph{Recall:} say the Latin for a wheel hub, then the Latin for a drag shoe, then say \emph{subiungō, subiungere} again
\end{itemize}

\begin{tcolorbox}[breakable,colback=teal!4,colframe=teal!35!black,title={Before you move on}]
What does subiungō, subiungere mean? (\textquotedblleft{}To yoke, to harness; to put under\textquotedblright{}.) Say it once more.
\end{tcolorbox}
\section[disiungō, disiungere]{disiungō, disiungere — to unyoke; to separate}

\label{lesson:LA-C430-disiungo}



\begin{quote}
Before the new one: say the Latin for a drag shoe, then the Latin for to yoke.
\end{quote}

\subsection*{You'll want to know: disiungō, disiungere}
\textbf{disiungō, disiungere} — \textquotedblleft{}to unyoke; to separate\textquotedblright{}.

\begin{quote}
Vesperī agricola bovēs disiungit et domum redit.
\end{quote}

\begin{quote}
\emph{In the evening the farmer unyokes the oxen and goes home.}
\end{quote}

\emph{disiungō, disiungere, disiūnxī, disiūnctum}. English \emph{disjoint} and \emph{disjunction} come from it.

\begin{cousinweb}
\emph{disiungō} is built from \emph{dis-} 'apart' and \emph{iungō} 'to join, to yoke' (EtymDB; Lewis \& Short). English \emph{disjoint} comes from it through French.
\end{cousinweb}

\begin{grammarlens}[title={What you've built}]
One more word for telling, explaining and describing at B1.
\end{grammarlens}

\subsection*{Your turn}
\begin{itemize}
\raggedright
  \item \emph{Say it:} \emph{disiungō, disiungere}
  \item \emph{Say it:} \emph{disiungō, disiungere}, once more
  \item \emph{Say it:} a sentence of your own with \emph{disiungō, disiungere}
  \item \emph{Recall:} say the Latin for a drag shoe, then the Latin for to yoke, then say \emph{disiungō, disiungere} again
\end{itemize}

\begin{tcolorbox}[breakable,colback=teal!4,colframe=teal!35!black,title={Before you move on}]
What does disiungō, disiungere mean? (\textquotedblleft{}To unyoke; to separate\textquotedblright{}.) Say it once more.
\end{tcolorbox}
\section[(dialogue)]{Roads, carts and ships}

\label{lesson:LA-R430-review-423-430}



\begin{quote}
Say the Latin for a footpath and a mountain track.
\end{quote}

\subsection*{Your turn}
Say these aloud:

\begin{itemize}
\raggedright
  \item a footpath, a mountain track, a side path, a slope, to step
  \item a crossroads, a fork in the road, a bend, a winding, to bend
  \item a raised roadbed, a flint, a raised edge, a wheel rut, to spread
  \item dusty, muddy, a footprint, weary, to wear out
  \item an overnight stop, a coachman, a post, travel money, to set out on
  \item a departure, a return, travel abroad, to travel abroad, to go forward
  \item a wagon, a light gig, a two-wheeled chariot, a covered carriage, a travelling coach
  \item a cart pole, a wheel hub, a drag shoe, to yoke, to unyoke
\end{itemize}

\begin{tcolorbox}[breakable,colback=teal!4,colframe=teal!35!black,title={Before you move on}]
A footpath? (\textbf{sēmita, sēmitae}.) A mountain track? (\textbf{callis, callis}.) A side path? (\textbf{trāmes, trāmitis}.) A slope? (\textbf{clīvus, clīvī}.) To step? (\textbf{gradior, gradī}.) A crossroads? (\textbf{compitum, compitī}.) A fork in the road? (\textbf{bivium, biviī}.) A bend? (\textbf{flexus, flexūs}.) A winding? (\textbf{anfrāctus, anfrāctūs}.) To bend? (\textbf{flectō, flectere}.) A raised roadbed? (\textbf{agger, aggeris}.) A flint? (\textbf{silex, silicis}.) A raised edge? (\textbf{crepīdō, crepīdinis}.) A wheel rut? (\textbf{orbita, orbitae}.) To spread? (\textbf{sternō, sternere}.) Dusty? (\textbf{pulverulentus, pulverulenta, pulverulentum}.) Muddy? (\textbf{lutulentus, lutulenta, lutulentum}.) A footprint? (\textbf{vestīgium, vestīgiī}.) Weary? (\textbf{lassus, lassa, lassum}.) To wear out? (\textbf{dēfatīgō, dēfatīgāre}.) An overnight stop? (\textbf{mānsiō, mānsiōnis}.) A coachman? (\textbf{raedārius, raedāriī}.) A post? (\textbf{statiō, statiōnis}.) Travel money? (\textbf{viāticum, viāticī}.) To set out on? (\textbf{ingredior, ingredī}.) A departure? (\textbf{profectiō, profectiōnis}.) A return? (\textbf{reditus, reditūs}.) Travel abroad? (\textbf{peregrīnātiō, peregrīnātiōnis}.) To travel abroad? (\textbf{peregrīnor, peregrīnārī}.) To go forward? (\textbf{prōgredior, prōgredī}.) A wagon? (\textbf{carrus, carrī}.) A light gig? (\textbf{cisium, cisiī}.) A two-wheeled chariot? (\textbf{essedum, essedī}.) A covered carriage? (\textbf{carpentum, carpentī}.) A travelling coach? (\textbf{carrūca, carrūcae}.) A cart pole? (\textbf{tēmō, tēmōnis}.) A wheel hub? (\textbf{modiolus, modiolī}.) A drag shoe? (\textbf{sufflāmen, sufflāminis}.) To yoke? (\textbf{subiungō, subiungere}.) To unyoke? (\textbf{disiungō, disiungere}.)
\end{tcolorbox}
